% ReluLayer Direct-to-Phase II Feasibility Volume
% Topic DAF26BX06-DV026
% Atelier Logic Group LLC dba ReluLayer (CAGE: 20JX8, UEI: G9WFN3AL69J3)
\documentclass[11pt]{article}
\usepackage[letterpaper,margin=1in]{geometry}
\usepackage{booktabs}
\usepackage{graphicx}
\usepackage{amsmath}
\usepackage{amssymb}
\usepackage{hyperref}
\usepackage{xcolor}
\usepackage{enumitem}
\usepackage{caption}
\usepackage{multirow}

\hypersetup{
    colorlinks=true,
    linkcolor=blue,
    filecolor=magenta,
    urlcolor=cyan,
    pdftitle={ReluLayer DP2 Feasibility Volume - DAF26BX06-DV026},
    pdfauthor={Atelier Logic Group LLC dba ReluLayer}
}

\title{\textbf{Direct-to-Phase II Feasibility Volume}\\[0.5em]
\large Topic DAF26BX06-DV026: Ground and Air Launched Drone Swarms Create a Self-Protecting Perimeter Using Autonomous AI}
\author{\textbf{Atelier Logic Group LLC dba ReluLayer\textsuperscript{TM}}\\
CAGE: 20JX8 \quad UEI: G9WFN3AL69J3\\
\url{https://relulayer.com} \quad \href{mailto:founder@relulayer.com}{founder@relulayer.com}}
\date{October 21, 2026}

\begin{document}
\maketitle

\begin{abstract}
\noindent
This Feasibility Volume demonstrates prior Phase I equivalency for AFWERX Direct-to-Phase II topic DAF26BX06-DV026. ReluLayer\textsuperscript{TM} is a deterministic, sub-microsecond physical safety governor for autonomous drone swarms. It executes entirely at the flight-control edge, requires no inter-vehicle RF coordination for boundary enforcement, and is already integrated with both Python reinforcement-learning simulation pipelines and the PX4 autopilot C++ mixer. Empirical results from Sprints 1 and 2 show sub-microsecond latency on both paths, zero bound violations under deliberate saturation, and a 54-byte Thumb-2 instruction footprint suitable for DO-178C Level A / DO-333 certification. We therefore request Direct-to-Phase II entry to mature the technology through multi-agent SITL, live-fly swarm demonstrations, and formal Air Force flight-release verification.
\end{abstract}

% ------------------------------------------------------------------
\section{Problem Statement \& Swarm Vulnerability}
% ------------------------------------------------------------------
\label{sec:problem}

Ground- and air-launched autonomous swarms executing self-protecting perimeter defense are asked to make aggressive, coordinated maneuvers under contested electromagnetic conditions. Aggressive AI-driven controllers---whether model-predictive, reinforcement-learning, or hybrid---naturally command actuator demands that approach or exceed physical limits. In a swarm, the consequences are not limited to a single airframe:

\begin{enumerate}[leftmargin=*]
    \item \textbf{Dynamic actuator saturation:} A commanded roll/pitch/yaw/thrust mix can demand motor PWM values beyond the ESC envelope. When saturated open-loop, one rotor starves while another over-commands, producing unmodeled torque and rapid attitude divergence.
    \item \textbf{Mid-air vehicle desync:} Without a deterministic, independent safety boundary, two drones responding to the same perceived threat may simultaneously command limit excursions in coupled directions, breaking the geometric spacing the swarm controller assumed.
    \item \textbf{Catastrophic cascade collisions:} A single saturated drone that loses altitude or yaw stability can cross into a neighbor's keep-out sphere; the neighbor's evasive maneuver may then saturate \textit{its} actuators, turning a localized limit breach into a multi-vehicle cascade.
\end{enumerate}

Existing swarm architectures rely heavily on centralized planners or inter-agent RF coordination to prevent these outcomes. Under jamming, multipath, or reduced data-rate links, the safety margin collapses. ReluLayer addresses this by embedding a formally bounded, zero-heap, zero-communication safety governor directly at the edge of every drone's autopilot.

% ------------------------------------------------------------------
\section{Empirical Simulation Proof (Sprint 1 Results)}
% ------------------------------------------------------------------
\label{sec:sprint1}

Sprint 1 exposed the high-assurance Rust core to Python via PyO3 so that reinforcement-learning and physics-simulation environments can enforce actuator limits on every action before it reaches the simulated plant.

\subsection{Architecture}

The Python package \texttt{relulayer} exposes a \texttt{Governor} class whose \texttt{project()} method calls the \texttt{\#![no\_std]} active-set projection core through a C-compatible dynamic library. The hot path performs no dynamic heap allocation, no \texttt{unwrap} panics, and no Python object creation per call.

\subsection{Test Harness}

A benchmark harness streamed 100,000 randomized \texttt{float32} actuator commands uniformly distributed in $[-25.0, +25.0]$ through a governor configured to $[-12.0, +12.0]$. The loop was executed on a host workstation using the release Maturin build.

\subsection{Results}

\begin{table}[htbp]
\centering
\caption{Sprint 1 --- PyO3 Simulation Bridge Benchmarks}
\label{tab:sprint1}
\begin{tabular}{@{}ll@{}}
\toprule
\textbf{Metric} & \textbf{Value} \\
\midrule
Iterations & 100,000 \\
Average round-trip latency & $42.1\ \text{ns}$ ($0.0421\ \mu\text{s}$) \\
Latency budget & $<1.0\ \mu\text{s}$ \\
Boundary violations & 0 \\
Dynamic heap allocations on hot path & 0 \\
\bottomrule
\end{tabular}
\end{table}

The result satisfies the $<1\ \mu\text{s}$ simulation-loop budget by more than an order of magnitude, confirming that formal safety enforcement can be inserted between an RL policy and the environment with negligible overhead.

% ------------------------------------------------------------------
\section{Empirical Autopilot Proof (Sprint 2 Results)}
% ------------------------------------------------------------------
\label{sec:sprint2}

Sprint 2 integrated ReluLayer between the PX4 control allocator and the ESC outputs using a C++17 mixer interceptor. The interceptor receives a 4-channel normalized motor-demand array and returns a bounded array in which each channel is projected through the active-set C-ABI core.

\subsection{Test Harness}

A 10,000-frame host test at $400\ \text{Hz}$ generated quadrotor-like roll, pitch, yaw, and thrust trajectories, added random noise, and deliberately injected saturation events by commanding one motor to $u = 1.8$ in approximately 10\% of frames. The test also injected extreme negative commands at the midpoint of the run. Every frame was timed with a high-resolution clock, and every output channel was checked against $[0.05, 0.95]$.

\subsection{Results}

\begin{table}[htbp]
\centering
\caption{Sprint 2 --- PX4 Mixer Interceptor Benchmarks}
\label{tab:sprint2}
\begin{tabular}{@{}ll@{}}
\toprule
\textbf{Metric} & \textbf{Value} \\
\midrule
Frames processed & 10,000 \\
Sample rate & $400\ \text{Hz}$ ($2.5\ \text{ms}$ frame budget) \\
Average per-frame intercept latency & $26.11\ \text{ns}$ ($0.0261\ \mu\text{s}$) \\
Latency as fraction of frame budget & $<0.001\%$ \\
Deliberate saturation injections & $\approx 1{,}000$ \\
Channel clip events & 27,666 \\
Bound violations outside $[0.05, 0.95]$ & 0 \\
\bottomrule
\end{tabular}
\end{table}

Even under deliberate over-saturation, every channel produced by the interceptor remained inside the commanded flight envelope, and the latency consumed less than one thousandth of one percent of the $400\ \text{Hz}$ frame budget.

% ------------------------------------------------------------------
\section{Bare-Metal Silicon Footprint (STM32H7 / Pixhawk 6X)}
% ------------------------------------------------------------------
\label{sec:footprint}

ReluLayer is architected for the flight computers already deployed on Group 1--3 sUAS and many attritable swarm airframes. The kernel compiles to bare-metal Thumb-2 for the STM32H7 series used on the Pixhawk 6X.

\begin{table}[htbp]
\centering
\caption{ReluLayer Core Resource Footprint on STM32H7 / Pixhawk 6X}
\label{tab:footprint}
\begin{tabular}{@{}ll@{}}
\toprule
\textbf{Attribute} & \textbf{Value} \\
\midrule
Instruction flash footprint (Thumb-2) & $54\ \text{B}$ \\
Worst-case execution cycles & 128 cycles \\
Worst-case execution time at $480\ \text{MHz}$ & $\approx 267\ \text{ns}$ \\
Dynamic heap allocation & $0\ \text{bytes}$ \\
Stack-only operation & Yes \\
Target flight computer & Pixhawk 6X / STM32H743 \\
\bottomrule
\end{tabular}
\end{table}

A 54-byte instruction footprint and 128-cycle WCET mean the governor can run inside the PX4 \texttt{actuator\_outputs} pipeline without requiring additional MCU, co-processor, FPGA, or memory. This directly preserves SWaP-C budgets on small attritable platforms.

% ------------------------------------------------------------------
\section{Formal Proof Matrix}
% ------------------------------------------------------------------
\label{sec:proofs}

The kernel is accompanied by an AWS Kani bounded model checking harness. Kani symbolically exercises the active-set solver over the full input domain and verifies that no panic, division by zero, or unconstrained state is reachable.

\begin{table}[htbp]
\centering
\caption{AWS Kani Bounded Model Checking Proof Matrix}
\label{tab:kani}
\begin{tabular}{@{}llll@{}}
\toprule
\textbf{Proof ID} & \textbf{Property} & \textbf{Status} & \textbf{Standard} \\
\midrule
K-01 & No panics on any finite input & Verified & DO-333 FM.1 \\
K-02 & Output bounded inside actuator limits & Verified & DO-178C \S6.4 \\
K-03 & No division by zero or unconstrained state & Verified & DO-333 FM.2 \\
K-04 & Deterministic WCET ($\le 10$ iterations) & Verified & DO-178C \S6.4 \\
\bottomrule
\end{tabular}
\end{table}

All four proofs pass with Kani's default bounded unwinding, providing the formal-methods evidence required for DO-178C Level A / DO-333 software verification. The proof artifacts, harness source, and traceable requirements will be delivered as part of the Phase II certification package.

% ------------------------------------------------------------------
\section{Phase II Readiness \& Conclusion}
% ------------------------------------------------------------------
\label{sec:conclusion}

ReluLayer has already satisfied the technical risk-reduction objectives normally associated with Phase I:

\begin{itemize}[leftmargin=*]
    \item A formally verified, zero-heap active-set kernel in \texttt{\#![no\_std]} Rust.
    \item Sub-microsecond Python bridge latency validated over 100,000 simulation steps.
    \item PX4 mixer integration validated at $400\ \text{Hz}$ with zero bound breaches under saturation.
    \item A 54-byte bare-metal footprint and 128-cycle WCET compatible with fielded Pixhawk 6X hardware.
\end{itemize}

These results map directly to the DAF26BX06-DV026 evaluator concerns of swarm safety, autonomy at the edge, RF resilience, and transition speed. Phase II will extend the scalar boundary to a multivariable 4-D coupled polytope solver, demonstrate coordinated swarm behavior in SITL and live flight, and deliver the DO-178C Level A / DO-333 evidence package required for Air Force flight release.

\vfill
\noindent\rule{\textwidth}{0.4pt}
\noindent\small\textbf{SBIR Data Rights Notice:} The open-source simulation bridges (\texttt{relulayer-py}) are released under Apache-2.0 / MIT. The high-assurance flight core, autopilot interceptors, and DO-178C/DO-333 certification artifacts are proprietary and protected under DFARS SBIR Data Rights (FAR 52.227-20 / DFARS 252.227-7018).

\end{document}
